\title{Voids in morphology}
\subtitle{Exploring “uninflectedness”}
\author{Enrique L. Palancar and Sebastian Fedden}
\BackBody{If we were to imagine an inflectional morphological system as a galactic universe with its matter and governing laws, ``uninflectedness'' would mark the outer boundaries of this system, the ``voids'', zones where the gravitational pull of inflection weakens and different structural principles come into play. Uninflectedness helps trace the limits of inflectional organization and offers valuable insights into how inflection operates and interacts with its periphery. Uninflectedness is thus a complex phenomenon that this book examines in depth for the first time, assembling a collection of 13 edited chapters that explore the phenomenon both theoretically and descriptively with the purpose of building a better typologically informed theory of inflection.}

\renewcommand{\lsSeries}{eotms}
\renewcommand{\lsSeriesNumber}{17}
\renewcommand{\lsID}{560}
\renewcommand{\lsISBNdigital}{978-3-96110-568-7}
\renewcommand{\lsISBNhardcover}{978-3-98554-189-8}
\BookDOI{10.5281/zenodo.19222401}
\proofreader{Matthew Korte}
\typesetter{Lucie Janků, Bruno Behling}




